\chapter{Long-Range Dependency as an Evaluation Target}

\section{From Distant Information to Demonstrated Dependence}

Some graph predictions require information that is not local to the prediction target.
Graph distance identifies where information lies, but it does not establish whether that information determines the correct output.
A long-range dependency is therefore a property of a label function rather than a property of graph size or distance alone.
The correct output must depend on relevant distant information and, in the path-aware setting studied here, on how that information is connected to the target.

Figure~\ref{fig:glora-running-example} provides a paired example that separates reachability, correlation, and dependence.
Suppose a source node \(v_0\) carries a designated feature and a target node \(v_d\) has label True only when every internal node on a source-to-target path carries a required continuation feature.
Compare this graph with an otherwise identical graph in which the feature of one internal node is changed and the target label is False.
The target remains reachable from the source through the same edges in both graphs, so reachability does not explain the label difference.
The source feature is also present in both graphs, so its presence does not distinguish the labels in this pair.
The change in the correct label is instead tied to the changed path information, which demonstrates dependence under the stated label function.
By contrast, frequent co-occurrence of the source feature and the True label in unpaired data would establish correlation only.

This distinction matters across graph domains in which distant evidence can become relevant through relations.
Examples include social influence, upstream traffic conditions, non-local biochemical interactions, web relations, and transaction chains in anti-money-laundering analysis \cite{zhou2025glora}.
The entities and edges differ across these domains, but the evaluation question is the same: whether the correct output requires information connected across several hops.

\section{Path-Aware Dependencies}

GLoRa gives one formalisation of this requirement through a path-aware dependency \cite{zhou2025glora}.
Let \(f\) be a node-classification function that takes a graph and a target node as input, and let \(G=(V,E,\lambda)\).
Let
\[
p=(v_0,\ldots,v_d)
\]
be a simple directed path of length \(d\).
The nodes on the path are distinct, and \((v_{i-1},v_i)\in E\) for every \(i\) with \(1\leq i\leq d\).
The node \(v_0\) is the source node, \(v_d\) is the target node, and each \(v_i\) with \(0<i<d\) is an internal node.

For every internal node \(v_i\), choose an alternative feature vector \(a_i\neq\lambda(v_i)\).
Let \(G_i\) be the graph obtained by replacing only \(\lambda(v_i)\) with \(a_i\).
For indices \(0<k\leq m<d\), let \(G_{p[k,m]}\) retain \(G\) and add a duplicate of each node from \(v_k\) to \(v_m\).
Each duplicate has the same feature as its original node.
For every edge incident to a node in the segment, \(G_{p[k,m]}\) also adds the corresponding edge with each endpoint in the segment replaced by its duplicate.
Write \(v'_j\) for the duplicate of \(v_j\).
Let \(G_{p[k,m],j}\) replace the feature of \(v'_j\) with \(a_j\), and let \(G_{p[k,m],i,j}\) additionally replace the feature of the original node \(v_i\) with \(a_i\).

The function \(f\) relies on a path-aware dependency of length \(d\) if there exist a graph, path, and alternative vectors as defined above for which the following three conditions hold.
\begin{enumerate}
\item \(f(G,v_d)\neq f(G_i,v_d)\) for every \(i\) with \(0<i<d\).
\item \(f(G,v_d)=f(G_{p[k,m],j},v_d)\) for every \(k,j,m\) with \(0<k\leq j\leq m<d\).
\item \(f(G,v_d)\neq f(G_{p[k,m],i,j},v_d)\) for every \(i,k,j,m\) with \(0<i<d\) and \(0<k\leq j\leq m<d\).
\end{enumerate}

The first condition makes the classification sensitive to the designated perturbation at every internal node on the witness path.
This sensitivity alone does not show that the function uses those nodes as a path because an internal feature could matter for an unrelated reason.
The second condition requires the classification to remain unchanged when the feature of any duplicated node \(v'_j\) in any allowed segment \(p[k,m]\) is replaced by \(a_j\), while the original path remains intact.
The third condition requires the classification to change when the features of any original internal node \(v_i\) and any duplicated node \(v'_j\) in any allowed segment are replaced by \(a_i\) and \(a_j\), respectively; \(i\) and \(j\) need not be equal.
Together, the duplicate-only and joint perturbations tie the classification to an intact source-to-target path rather than to a detached feature test.

This path-aware definition is one formalisation of a long-range dependency rather than a claim that every long-range graph relation has this form.
Other definitions could describe distant-node interactions without a single target node or without making one path the witness.
The path-aware choice is appropriate here because it gives a precise task-level counterpart to the local message-passing computation introduced in Chapter~2.

\section{Distance and Dependency Length}

Graph distance, path length, dependency length, and path-aware dependency name different objects.
For nodes \(u\) and \(v\), the graph distance \(\operatorname{dist}_G(u,v)\) is the length of a shortest directed path from \(u\) to \(v\), when such a path exists.
Graph distance is a structural property of \(G\) and does not involve a label function.

Path length belongs to one particular path.
The path \(p=(v_0,\ldots,v_d)\) has length \(d\) because it contains exactly \(d\) edges, whether or not it is a shortest path between its endpoints.
Consequently, \(\operatorname{dist}_G(v_0,v_d)\) can be smaller than the length of \(p\).

Dependency length belongs to the task.
In the formalisation above, \(d\) is the exact length of a path that witnesses the required dependence of the label function.
Path-aware dependency names this dependence relation, whereas dependency length names the number of edges on its witness path.
Distance alone is not a dependency because it says nothing about whether changing distant information can change the output.

Graph distance, path length, and dependency length can coincide numerically, but any equality between them must be established.
A graph can contain a long path while its labels depend only on local information.
A graph can also have a large diameter without any label depending on a pair of distant nodes.
A chosen witness path can also be longer than the shortest route between its endpoints.
An intended rule that refers to a path of length \(d\) is therefore insufficient if the observed labels can be fitted without using that path.

The exact witness-path length used in this definition must also remain distinct from the length parameter supplied to the GLoRa generator in Chapter~6.
A generator parameter controls a construction, so its relation to the exact length of a witness path in a generated graph must be stated separately rather than assumed.
Within this chapter, \(d\) denotes the exact length of the witness path unless a different role is stated explicitly.

Once dependency length has been defined as a task requirement, an evaluation can vary it under controlled conditions.
This replaces a single score with a performance profile over specified witness-path lengths.
Such a profile describes how far demonstrated use extends under the evaluation conditions without treating distance alone as dependence.

\section{Scope of the Definition}

The three-condition definition is node-level.
Paper I applies the definition to GLoRa's inductive binary node-classification setting, and the paper's formal benchmark guarantee is limited to that setting \cite{zhou2025glora}.
The function \(f(G,v_d)\) assigns the target node one of the labels True and False.
The inductive setting requires the learned rule to apply to new graph instances rather than to memorised node identities in one fixed graph.
Within this thesis, formal claims based on the definition are limited to this setting unless an extension is stated explicitly.

The explicit target node makes the intervention direct.
Changing an internal node tests whether the output at \(v_d\) changes, while the duplicate-segment conditions test whether the output depends on an intact path to that target.
The source, internal nodes, and target therefore have fixed roles in the witness.

Graph-level tasks have no distinguished target node in their output format, but their labels can still depend on path information.
One extension treats the place where path information is assembled before graph readout as a conceptual target.
Another treats the path itself as the witness and asks whether perturbing its internal nodes changes the graph label.
Both extensions preserve the distinction between distant information being present and the output depending on its relation across the graph.

These graph-level extensions are conceptual and are not covered by Paper I's formal node-classification result.
A graph-level label may depend on several source nodes, several readout locations, or redundant witness paths.
The difficulty of learning such a label can therefore depend on the amount and arrangement of relevant information as well as on dependency length.

\section{Access Is Not Demonstrated Use}

Access is an architectural property relative to an input graph.
For a standard message-passing graph neural network, each layer expands the target node's receptive field by one directed hop \cite{scarselli2008graph,kipf2017semi}.
If a model has fewer than \(\operatorname{dist}_G(v_0,v_d)\) layers, the feature of \(v_0\) cannot affect the target representation through ordinary one-hop message passing.
Enough layers remove this structural barrier, but they establish only that the information can reach the target representation.

Demonstrated use is a claim about the learned function.
The learned prediction must respond to the relevant path interventions rather than merely receive the source feature inside its receptive field.
In the paired example above, a system that gives the same prediction after the internal feature changes has access to the source but does not exhibit the required dependence.
The duplicate-only and joint perturbations make this test stronger by checking that the response follows the intact path.

Correlation is weaker than both task dependence and demonstrated use.
A feature can predict the label because it co-varies with the label across the sampled examples, even when changing the relevant path information would leave the learned prediction unchanged.
High accuracy on such examples establishes predictive performance, but it does not by itself identify which information the learned function used.
Paired or intervention-based evidence is needed to separate the intended dependence from label-correlated alternatives.

The same boundary applies to failure.
A decline in accuracy as dependency length increases is a behavioural symptom, not a diagnosis of the mechanism that caused the decline.
The mechanism requires separate internal evidence with a prediction that distinguishes it from other explanations.
Chapter~4 develops that diagnostic standard after the learning target has been fixed.

\section{Evaluation Contract}

An evaluation supports a claim about learning a path-aware dependency of a specified length only when four requirements hold.
First, it must state the target output, the path-aware dependency to be learned, and the exact witness-path length.
Second, the evaluated system must have architectural access to the relevant information, and the target function must be expressible by the evaluated model family \cite{barcelo2020logical,zhou2025glora}.
Third, the evidence must make successful prediction depend on the relevant path changes while holding reachability and label-correlated alternatives fixed.
Fourth, comparisons across dependency lengths and systems must use controlled conditions and must report any distinction between an exact witness-path length and a generator parameter.

Under this contract, successful prediction supports a bounded claim that the learned function uses the specified path-aware dependency in the evaluated setting.
Failure supports only the conclusion that, under those conditions, the evaluation did not show that the system had learned a function with the specified path-aware dependency.
Neither result identifies a failure mechanism from accuracy alone, and neither establishes capability for every form of long-range graph relation.

These four requirements provide the basis for RQ~1.
With the learning target and required evidence fixed, Chapter~4 develops mechanism-specific predicted traces for distinguishing explanations of failure.
Chapter~5 then turns this contract into benchmark criteria, including the control of alternative fitting functions and fair comparison between graph learning systems.
